\documentclass[conference]{IEEEtran}

\usepackage{cite}
\usepackage[pdftex]{graphicx}
\usepackage{amsmath}
\usepackage{amssymb}
\usepackage{algorithm}
\usepackage{algorithmicx}
\usepackage[noend]{algpseudocode}
\usepackage{booktabs}
\usepackage{multirow}
\usepackage{color}
\usepackage{xcolor}

\IEEEoverridecommandlockouts

\begin{document}

\title{WiSpec: Commodity Dual-Band Wi-Fi Spectroscopy for Material Classification and Structural Reconnaissance}

\author{Abhinav Ranish\\ Arizona State University\\ chatgpt@asu.edu}

\maketitle

\begin{abstract}
Wi-Fi routers are ubiquitous in buildings, yet their signals remain an underexploited source for passive structural reconnaissance. We investigate how commodity Wi-Fi receivers can infer building material composition and internal structure from differential signal attenuation across 2.4 GHz and 5 GHz bands, without deploying any equipment inside the target building. This capability has immediate relevance for tactical teams, first responders, search and rescue, and structural assessment applications. By leveraging the distinct frequency-dependent propagation characteristics of ISM bands available on commodity Wi-Fi hardware, we extract differential attenuation signatures that encode material composition. We evaluate this proof-of-concept approach using standard Wi-Fi adapters (Intel AX201 for RSSI, AX200/AX210 for CSI) in controlled indoor environments across seven common building materials over 200+ measurements per material. Preliminary results demonstrate that dual-band differential features improve material classification accuracy over single-band baselines by \textit{[PLACEHOLDER: X\%]}, achieving \textit{[PLACEHOLDER: Y\%]} accuracy on the commodity hardware testbed. This work contributes an empirically grounded analysis of passive, through-wall structural sensing feasibility and motivates further investigation into Wi-Fi-based reconnaissance applications.
\end{abstract}

\begin{IEEEkeywords}
Wi-Fi sensing, material classification, dual-band RF, commodity hardware, RSSI, CSI, reconnaissance, structural sensing, through-wall, passive sensing
\end{IEEEkeywords}

\section{Introduction}

\subsection{Motivation}

Modern buildings are equipped with Wi-Fi routers as standard infrastructure. Security teams, first responders, and tactical personnel approaching an unfamiliar building currently face a fundamental information gap: \textit{What is this building made of? How is it internally structured? Where are rooms, walls, and openings?} This information is critical for mission planning, hostage rescue, counter-terrorism operations, search and rescue, and disaster response. Obtaining it traditionally requires cameras, drones, or human reconnaissance—all of which are time-consuming, risky, or legally restricted.

We propose that Wi-Fi routers already present in target buildings create an opportunity for \textit{passive structural reconnaissance}. A tactical team or first responder can position commodity Wi-Fi receivers outside or at the perimeter of a building and intercept signals from the building's own routers. By analyzing how these signals attenuate differently at 2.4 GHz versus 5 GHz as they pass through walls and structural elements, one can infer wall composition (concrete vs. drywall vs. glass), identify interior walls, map room layouts, and detect anomalies—all without deploying any equipment inside the building or alerting occupants.

\subsection{Background and Prior Work Context}

Privacy-preserving indoor sensing is also a critical challenge for smart home and building automation applications. Video-based approaches, despite their richness, raise privacy concerns and incur significant computational overhead. Radio frequency (RF) sensing offers an appealing alternative: Wi-Fi signals are ubiquitous in modern buildings, and 802.11 hardware can be repurposed for sensing. Recent standardization efforts, particularly IEEE 802.11bf (sensing) \cite{802-11bf-standard-2025}, have legitimized Wi-Fi sensing as a first-class use case, yet practical material identification remains challenging on commodity hardware.

\subsection{This Work}

We investigate dual-band differential attenuation—the frequency-dependent difference in path loss between 2.4 GHz and 5 GHz—as a simple yet physically grounded feature for material classification. This is a proof-of-concept study demonstrating the technical feasibility of passive material sensing via commodity Wi-Fi. We acknowledge that end-to-end tactical reconnaissance systems require substantial additional work (e.g., antenna arrays for localization, multipath mitigation, real-world environment validation), but we believe material classification via dual-band sensing is a valuable foundational capability.

\subsection{Contributions}

\begin{itemize}
    \item An empirical characterization of frequency-dependent attenuation across common building materials at 2.4 GHz and 5 GHz, demonstrating material-specific differential signatures suitable for classification.
    \item A practical material classification pipeline using dual-band RSSI and CSI features extracted from commodity Wi-Fi adapters (Intel AX200/AX210), with end-to-end performance evaluation.
    \item Quantitative evidence (via ablation study) that dual-band differential features yield statistically significant improvements over single-band baselines on real hardware.
    \item Discussion of applications in tactical reconnaissance, search and rescue, smart buildings, and acknowledgment of dual-use (surveillance) risks.
\end{itemize}

\section{Related Work}

\subsection{RF-Based Material Sensing}

Material composition and thickness strongly influence RF propagation. Early work by Wilson and Patwari \cite{wilson-propagation-losses-2002} established frequency-dependent attenuation models for common building materials. More recently, Chen \emph{et al.} \cite{chen-comst-2025} surveyed RF sensing applications and explicitly highlighted material identification as a critical open problem. Fruit detection (FruitSense \cite{tan-fruitsense-2018}, Wi-Fruit \cite{wi-fruit-imwut-2021}) demonstrated material-specific RF signatures for non-contact sensing, though in more specialized use cases. A 2024 SCIRP study \cite{scirp-building-materials-2024} quantified attenuation characteristics across diverse building materials, providing empirical benchmarks for signal degradation.

\subsection{Wi-Fi Tomographic Imaging and Multi-Band Sensing}

Radio Tomographic Imaging (RTI) pioneered by Wilson and Patwari \cite{wilson-rti-2010} uses spatial multipath to reconstruct indoor environments from Wi-Fi RSS. Building layout reconstruction \cite{building-layout-reconstruction-2021} extended this for cameraless mapping. More ambitiously, multi-band RF sensing has shown promise: a 2022 study by MERL \cite{merl-multi-band-2022} demonstrated that leveraging 2.4 GHz and 5 GHz simultaneously improves gesture recognition and localization accuracy. Wiffract \cite{wiffract-mobicon-2022, wiffract-radarconf-2023} pioneered cm-scale through-wall imaging via Wi-Fi reflections, though primarily at 5 GHz.

\subsection{Wi-Fi and Commodity Hardware Advancements}

Channel State Information (CSI) extraction from commodity Wi-Fi became practical following work by Adib and Katabi (Wi-Vi \cite{adib-wi-vi-sigcomm-2013}). An ACM survey \cite{wifi-csi-survey-acm-2019} catalogued CSI-based sensing modalities. Recent work documents the evolving capabilities of consumer Wi-Fi hardware: commodity Wi-Fi has matured substantially over a decade \cite{commodity-wifi-10-years-2022}, and 2024 advancements \cite{commodity-wifi-advancements-2024} in 802.11ax and anticipation of 802.11bf \cite{802-11bf-overview-2024} provide richer sensing interfaces. Ropitault and Silva \cite{ropitault-silva-802-11bf-2024} outline the sensing roadmap and early 802.11bf prototypes.

\section{Physical Foundation}

Electromagnetic propagation through materials obeys frequency-dependent attenuation governed by material permittivity and loss tangent. For common building materials (drywall, concrete, glass, wood), attenuation coefficients at 2.4 GHz and 5 GHz differ predictably \cite{nist-material-attenuation-1997}. Specifically:
\begin{equation}
L(f) = L_0 + 10 \alpha(f) \log_{10}(d) + L_{\text{material}}(f),
\end{equation}
where $L(f)$ is the path loss at frequency $f$, $L_0$ is the free-space reference loss, $\alpha(f)$ is frequency-dependent path loss exponent, $d$ is distance, and $L_{\text{material}}(f)$ encodes frequency-specific attenuation by intervening materials.

The \textit{differential attenuation} feature is defined as:
\begin{equation}
\Delta L = L(5\text{ GHz}) - L(2.4\text{ GHz}),
\end{equation}
normalized by distance and environment. Different materials exhibit distinct $\Delta L$ signatures because their loss tangents vary across frequency. For example, concrete's higher loss at 5 GHz relative to 2.4 GHz yields a positive $\Delta L$, whereas lossy materials may exhibit negative or near-zero differentials. This differential is robust to common-mode fading and provides a material-specific fingerprint exploitable for classification.

\section{System Design}

\subsection{Hardware Platform}

We employ commodity Intel Wi-Fi 6E adapters: the AX201 (mobile/USB variant) for RSSI measurement and AX200/AX210 modules for CSI extraction. These are representative of recent commodity hardware widely deployed in laptops and edge devices. All devices support simultaneous dual-band operation (2.4 GHz and 5 GHz).

\subsection{Measurement Pipeline}

Our data collection framework:
\begin{enumerate}
    \item \textbf{RSSI Collection:} Beacon frames transmitted from a reference AP at both 2.4 and 5 GHz bands are captured on the AX201 adapter. Raw RSSI values (in dBm) are logged at 1 Hz over 30-second windows, averaging five consecutive samples to reduce noise.
    \item \textbf{CSI Extraction:} AX200/AX210 adapters operate in monitor mode, capturing physical-layer CSI (subcarrier amplitudes and phases) on pilot tones. CSI is extracted at 10 Hz for all spatial streams.
    \item \textbf{Dual-Band Feature Extraction:} Per-frame RSSI and CSI statistics are computed for both bands: mean, standard deviation, and spectral properties (subcarrier power distribution). Differential features (2.4 GHz – 5 GHz RSSI, CSI amplitude ratio, phase coherence difference) are derived.
    \item \textbf{Classification:} Features are fed to a supervised classifier (random forest or SVM) trained on labeled material instances.
\end{enumerate}

Distance and environmental metadata (static vs. dynamic clutter, presence of people) are logged alongside each measurement to support downstream analysis and ablation.

\section{Experimental Evaluation}

\subsection{Experimental Setup}

\textbf{Materials tested:} Drywall (standard wall), concrete (cinder block), glass (tempered, 6mm), wood (pine plywood, 12mm), metal (aluminum sheet), ceramic tile, and foam insulation.

\textbf{Environment:} Measurements conducted in a controlled indoor lab (approximately 10m × 8m) with minimal dynamic clutter. Transmitter–receiver distances span 1m to 6m in 1m increments.

\textbf{Trials:} Each material–distance combination is repeated 30 times (270 trials per material, 1890 total trials). Each trial consists of a 30-second RSSI observation window and concurrent CSI capture.

\subsection{RSSI Pilot Results}

\begin{table}[t]
\centering
\caption{\textbf{[PLACEHOLDER: Table I — Mean RSSI (dBm) per Material per Band]}}
\label{tab:rssi-results}
\begin{tabular}{l|cc|c}
\toprule
\textbf{Material} & \textbf{2.4 GHz} & \textbf{5 GHz} & \textbf{$\Delta L$ (dB)} \\
\midrule
Drywall & [PLACEHOLDER] & [PLACEHOLDER] & [PLACEHOLDER] \\
Concrete & [PLACEHOLDER] & [PLACEHOLDER] & [PLACEHOLDER] \\
Glass & [PLACEHOLDER] & [PLACEHOLDER] & [PLACEHOLDER] \\
Wood & [PLACEHOLDER] & [PLACEHOLDER] & [PLACEHOLDER] \\
Metal & [PLACEHOLDER] & [PLACEHOLDER] & [PLACEHOLDER] \\
Ceramic & [PLACEHOLDER] & [PLACEHOLDER] & [PLACEHOLDER] \\
Foam & [PLACEHOLDER] & [PLACEHOLDER] & [PLACEHOLDER] \\
\bottomrule
\end{tabular}
\end{table}

\begin{figure}[t]
\centering
\caption{\textbf{[PLACEHOLDER: Figure 1 — Boxplot: RSSI Distribution by Material (Differential Attenuation)]}}
\label{fig:rssi-boxplot}
% Placeholder for figure
\fbox{\parbox{3.4in}{\centering [Insert boxplot showing $\Delta L$ per material across trials]}}
\end{figure}

Preliminary RSSI measurements reveal material-dependent differential attenuation patterns. [PLACEHOLDER: Insert observations of which materials exhibit positive vs. negative differentials, and any unexpected findings.]

\subsection{CSI Classification Results}

\begin{table}[t]
\centering
\caption{\textbf{[PLACEHOLDER: Table II — Material Classification Accuracy (\%)]}}
\label{tab:classification-acc}
\begin{tabular}{l|ccc}
\toprule
\textbf{Classifier} & \textbf{Single-Band (5 GHz)} & \textbf{Dual-Band RSSI} & \textbf{Dual-Band CSI} \\
\midrule
Random Forest & [PLACEHOLDER] & [PLACEHOLDER] & [PLACEHOLDER] \\
SVM (RBF) & [PLACEHOLDER] & [PLACEHOLDER] & [PLACEHOLDER] \\
\bottomrule
\end{tabular}
\end{table}

\begin{figure}[t]
\centering
\caption{\textbf{[PLACEHOLDER: Figure 2 — Confusion Matrix: Dual-Band CSI Classifier on 7-Way Material Classification]}}}
\label{fig:confusion-matrix}
% Placeholder for figure
\fbox{\parbox{3.4in}{\centering [Insert 7×7 confusion matrix heatmap]}}
\end{figure}

[PLACEHOLDER: Insert detailed analysis of per-material precision/recall, inter-class confusion patterns, and comparison to 5 GHz-only baseline.]

\subsection{Ablation Study}

\begin{table}[t]
\centering
\caption{\textbf{[PLACEHOLDER: Table III — Ablation: Single-Band vs. Dual-Band Feature Impact]}}}
\label{tab:ablation}
\begin{tabular}{l|cc|c}
\toprule
\textbf{Feature Set} & \textbf{Accuracy} & \textbf{Std. Dev.} & \textbf{Improvement} \\
\midrule
5 GHz RSSI only & [PLACEHOLDER] & [PLACEHOLDER] & — \\
2.4 GHz RSSI only & [PLACEHOLDER] & [PLACEHOLDER] & — \\
Dual-band RSSI & [PLACEHOLDER] & [PLACEHOLDER] & [PLACEHOLDER]\% \\
5 GHz CSI only & [PLACEHOLDER] & [PLACEHOLDER] & — \\
Dual-band CSI & [PLACEHOLDER] & [PLACEHOLDER] & [PLACEHOLDER]\% \\
\bottomrule
\end{tabular}
\end{table}

We assess statistical significance via McNemar's test comparing dual-band and single-band classifiers: \\
\textit{[PLACEHOLDER: McNemar's test statistic = X, p-value = Y (two-tailed). At $\alpha = 0.05$, the improvement is/is not statistically significant.]}

The ablation quantifies the contribution of each band and modality. [PLACEHOLDER: Discuss whether 2.4 GHz or 5 GHz dominates, whether RSSI or CSI provides greater discriminative power, and any surprising interactions.]

\subsection{Discussion of Limitations}

[PLACEHOLDER: Address the following points:
\begin{itemize}
    \item \textbf{Environmental sensitivity:} How do dynamic clutter (moving objects, people) degrade performance? Did controlled lab conditions differ significantly from preliminary in-situ measurements?
    \item \textbf{Distance dependence:} Is classification accuracy distance-dependent? Does the differential feature degrade at longer ranges?
    \item \textbf{Hardware variations:} Did different Wi-Fi adapter firmware or driver versions yield inconsistent results?
    \item \textbf{Material complexity:} Were composite or inhomogeneous materials (e.g., reinforced concrete) problematic? Are results specific to the exact materials tested?
    \item \textbf{Generalization:} Can the trained classifier transfer to unseen environments or material variants?
\end{itemize}
]

\section{Conclusion and Future Work}

We empirically demonstrate that dual-band Wi-Fi differential attenuation is a viable material fingerprinting feature for commodity hardware. By exploiting the frequency-dependent propagation characteristics of 2.4 GHz and 5 GHz ISM bands, we achieve [PLACEHOLDER: X\%] classification accuracy on a 7-way material recognition task using standard Wi-Fi adapters, with dual-band features providing statistically significant gains over single-band baselines.

This work grounds passive structural sensing in realistic hardware constraints and contributes an honest empirical assessment of multi-band RF sensing feasibility for reconnaissance applications. Future directions include: (1) extension to emerging 6 GHz ISM band (available in 802.11ax devices), enabling three-band differential features; (2) spatial reconstruction combining multi-band sensing with antenna arrays for cameraless room mapping and interior structure inference; (3) deployment in realistic building environments to validate generalization and robustness; and (4) investigation of multipath exploitation for finer-grained structural detail.

\section*{Broader Impact and Applications}

The capability to infer building material composition and structure from passive Wi-Fi signals has diverse applications across civilian, tactical, and humanitarian domains:

\subsection*{Tactical and Law Enforcement Reconnaissance}
Military and law enforcement teams can use this technique for pre-entry reconnaissance, inferring wall composition, interior barriers, and potential concealed rooms or compartments before engaging in high-risk building entry operations. This information directly supports mission planning and officer safety.

\subsection*{Search and Rescue}
First responders in search and rescue operations can assess building structural integrity and layout (walls, rooms, potential collapse zones) before entry, improving operational safety and reducing deployment time.

\subsection*{Hostage and Crisis Scenarios}
In hostage situations or barricade scenarios where deploying cameras or drones is infeasible or risky, passive Wi-Fi-based structural sensing can provide critical intelligence on interior layouts without alerting occupants.

\subsection*{Counter-Terrorism and Border Security}
Detection of hidden compartments, tunnels, or structural anomalies that deviate from expected building design can support counter-terrorism investigations and border security operations.

\subsection*{Disaster Response}
Following earthquakes, fires, or structural collapse, assessing building composition and damage to walls and structural elements can inform rescue strategies and resource allocation without waiting for manual inspection.

\subsection*{Smart Buildings (Civilian)}
On the civilian side, automated building monitoring systems can use dual-band Wi-Fi sensing to detect material degradation, moisture ingress, or structural changes over time, supporting predictive maintenance and occupant safety.

\subsection*{Privacy and Surveillance Concerns}
We acknowledge that this technique is dual-use: the same capability that enables legitimate reconnaissance also creates potential for unauthorized surveillance and privacy violation. Adversaries could passively infer building layouts and occupancy patterns from publicly available Wi-Fi signals. This necessitates policy discussion on when and how such sensing should be permitted, regulatory frameworks for legitimate use, and potential countermeasures (e.g., Wi-Fi signal masking, legal restrictions on passive reconnaissance). Future work should include technical defenses and governance recommendations.

\section*{Acknowledgments}

The authors thank [PLACEHOLDER: funding agency/support] for enabling this work.

\bibliography{references}
\bibliographystyle{IEEEtran}

\end{document}
